\documentclass[11pt]{article}
\usepackage[margin=1in]{geometry}
\usepackage{hyperref}
\usepackage{graphicx}
\usepackage{enumitem}
\usepackage{array}
\usepackage{longtable}
\usepackage{amsmath}

\hypersetup{
  colorlinks=true,
  linkcolor=blue,
  urlcolor=blue
}

\setlength{\parindent}{0pt}
\setlength{\parskip}{6pt}
\setcounter{tocdepth}{2}

\title{COS Open-Book Notes (Lectures 10--20)}
\author{Ashoka University COS}\date{April 23, 2026}

\begin{document}
\maketitle

\tableofcontents
\newpage

\section*{How to Use This Document to Pass the Exam}
\addcontentsline{toc}{section}{How to Use This Document to Pass the Exam}
\begin{itemize}[leftmargin=*]
  \item The \textbf{Index} (Table of Contents) lists every concept with real page numbers after PDF compilation. Use it as your entry point.
  \item For each concept, read the \textbf{Step-by-step} list first, then the \textbf{Examples}. This mirrors how exam questions are structured.
  \item Use the \textbf{Common Checks and Traps} line to avoid typical mistakes under time pressure.
  \item If you are stuck, jump to the \textbf{MIPS Programming Guide} appendix for syntax, syscall, and pattern refreshers.
  \item When you get additional past papers, we can replace or tune examples to match difficulty and style.
\end{itemize}

\section*{Diagram Placeholders}
\addcontentsline{toc}{section}{Diagram Placeholders}
These notes include placeholders for diagrams. Provide screenshots and place them in the \texttt{images/} folder as requested. I will wire them into the document.

\newpage

\section{Floating Point and Computer Arithmetic (Lectures 10--11)}

\subsection{IEEE 754 Formats and Fields (L10, L11)}
\textbf{What it is.} Floating point numbers are stored as sign, exponent, and fraction. Single precision: 1 sign bit, 8 exponent bits, 23 fraction bits. Double precision: 1 sign bit, 11 exponent bits, 52 fraction bits. Normalized value:
\[\text{value} = (-1)^S \times (1.F) \times 2^{E-\text{bias}}\]
where bias is 127 (single) or 1023 (double).

\textbf{Step-by-step.}
\begin{enumerate}[leftmargin=*]
  \item Split the bits into S, E, and F.
  \item Restore the hidden 1 for normalized values (E not all zeros or ones).
  \item Compute $e = E - \text{bias}$.
  \item Compute $(1.F)$ as a binary fraction and apply sign.
\end{enumerate}

\textbf{Examples.}
\begin{enumerate}[leftmargin=*]
  \item Single: $S=0, E=127, F=0$ gives $+1.0$.
  \item Single: $S=1, E=128, F=1000...$ gives $-1.5 \times 2^1 = -3.0$.
  \item Double: $S=0, E=1023, F=0$ gives $+1.0$.
\end{enumerate}

\textbf{Common checks and traps.} Do not forget the hidden 1 for normalized values.

\begin{figure}[h]
\centering
\fbox{\parbox{0.9\linewidth}{Diagram placeholder: IEEE 754 single and double field layout (L10)}}
\end{figure}

\subsection{Exponent Bias, Encoding, and Special Exponents (L10, L11)}
\textbf{What it is.} Exponents are stored with a bias so they can be compared as unsigned. Exponent all-zeros and all-ones are reserved for special values.

\textbf{Step-by-step.}
\begin{enumerate}[leftmargin=*]
  \item Convert actual exponent $e$ to stored exponent $E = e + \text{bias}$.
  \item If $E=0$, the number is zero or denormal.
  \item If $E$ is all ones, the number is infinity or NaN.
\end{enumerate}

\textbf{Examples.}
\begin{enumerate}[leftmargin=*]
  \item Single: $e=-3 \Rightarrow E=124$ (binary 01111100).
  \item Single: $e=0 \Rightarrow E=127$ (binary 01111111).
  \item Single: $E=255$ implies infinity or NaN (check fraction).
\end{enumerate}

\textbf{Common checks and traps.} Do not treat $E=0$ as a normal exponent.

\subsection{Normalized, Denormal, Zero, Infinity, and NaN (L10)}
\textbf{What it is.} Normalized values use a hidden 1. Denormals use exponent $E=0$ and hidden bit 0 for gradual underflow. Special values use $E$ all ones.

\textbf{Step-by-step.}
\begin{enumerate}[leftmargin=*]
  \item Normalized: $1.F \times 2^{E-\text{bias}}$.
  \item Denormal: $0.F \times 2^{1-\text{bias}}$ (single: $2^{-126}$).
  \item Zero: $E=0$ and $F=0$ (both +0 and -0).
  \item Infinity: $E$ all ones and $F=0$.
  \item NaN: $E$ all ones and $F \ne 0$.
\end{enumerate}

\textbf{Examples.}
\begin{enumerate}[leftmargin=*]
  \item Smallest positive normal (single): $1.0 \times 2^{-126}$.
  \item Smallest positive denormal (single): $2^{-149}$.
  \item Overflow in a multiplication produces $+\infty$.
\end{enumerate}

\textbf{Common checks and traps.} Denormals have lower precision and a hidden 0, not 1.

\begin{figure}[h]
\centering
\fbox{\parbox{0.9\linewidth}{Diagram placeholder: Normalized and denormal range on real line (L10)}}
\end{figure}

\subsection{Encoding and Decoding Procedure (L11)}
\textbf{What it is.} Steps to map between decimal values and IEEE 754 bit patterns.

\textbf{Step-by-step (encode).}
\begin{enumerate}[leftmargin=*]
  \item Determine sign $S$.
  \item Convert magnitude to binary and normalize as $1.F \times 2^e$.
  \item Compute $E=e+\text{bias}$ and fill fraction bits.
  \item Round to fit the fraction width.
\end{enumerate}

\textbf{Step-by-step (decode).}
\begin{enumerate}[leftmargin=*]
  \item Read $S$, $E$, $F$ and detect special values.
  \item Restore the hidden bit (if normalized).
  \item Compute the final value using $(-1)^S (1.F) 2^{E-\text{bias}}$.
\end{enumerate}

\textbf{Examples.}
\begin{enumerate}[leftmargin=*]
  \item Encode $-0.75$: $-0.11_2 = -1.1_2 \times 2^{-1}$, $E=126$, fraction $1000...$.
  \item Decode $1\ 01111110\ 1100...$: sign negative, $E=126$, value $-0.875$.
  \item Encode $6.5$: $110.1_2 = 1.101_2 \times 2^2$, $E=129$, fraction $1010...$.
\end{enumerate}

\textbf{Common checks and traps.} Normalization changes exponent; do not drop it.

\subsection{Floating-Point Addition and Subtraction Algorithm (L10, L11)}
\textbf{What it is.} FP add/sub aligns exponents, adds significands, normalizes, and rounds. Signs are handled using 2's complement when needed.

\textbf{Step-by-step.}
\begin{enumerate}[leftmargin=*]
  \item Align binary points by shifting the smaller-exponent significand right.
  \item Add or subtract significands (use 2's complement for opposite signs).
  \item Normalize the result (shift left or right; adjust exponent).
  \item Round to target precision; renormalize if rounding carries.
  \item Check for overflow/underflow.
\end{enumerate}

\textbf{Examples.}
\begin{enumerate}[leftmargin=*]
  \item $1.111_2 \times 2^{-1} + 1.011_2 \times 2^{-3}$: align, add, normalize, round.
  \item $1.000_2 \times 2^{-3} + (-1.000_2) \times 2^{2}$: convert to 2's complement, add, normalize.
  \item Opposite signs with close magnitudes: result needs leading-zero normalization.
\end{enumerate}

\textbf{Common checks and traps.} Alignment can drop low bits; use guard/round/sticky.

\subsection{Floating-Point Multiplication Algorithm (L10)}
\textbf{What it is.} Multiply significands, add exponents (minus bias), compute sign, normalize, round, and check range.

\textbf{Step-by-step.}
\begin{enumerate}[leftmargin=*]
  \item Sign: $S = S_X \oplus S_Y$.
  \item Exponent: $E_Z = E_X + E_Y - \text{bias}$.
  \item Multiply significands (may be 2p bits).
  \item Normalize (shift right at most 1; adjust exponent).
  \item Round and check overflow/underflow.
\end{enumerate}

\textbf{Examples.}
\begin{enumerate}[leftmargin=*]
  \item $(1.010)_2 \times 2^{-1} \times (-1.110)_2 \times 2^{-2}$: sign negative, exponent $-3$.
  \item Multiply by $2^k$: only exponent changes (fast path).
  \item Product of two numbers just under 2.0 may require normalization shift.
\end{enumerate}

\textbf{Common checks and traps.} The product can have 2 leading bits; normalize once.

\subsection{Guard, Round, and Sticky Bits (L10)}
\textbf{What it is.} Extra bits retained during alignment and arithmetic to improve rounding accuracy.

\textbf{Step-by-step.}
\begin{enumerate}[leftmargin=*]
  \item Guard bit: first bit beyond fraction.
  \item Round bit: next bit after guard.
  \item Sticky bit: OR of all remaining shifted-out bits.
  \item Use these bits to decide round-to-nearest (or other mode).
\end{enumerate}

\textbf{Examples.}
\begin{enumerate}[leftmargin=*]
  \item Guard=1, sticky=1 usually rounds up.
  \item Guard=0, sticky=1 usually rounds down.
  \item Guard=1, round=0, sticky=0 is a tie (round-to-even).
\end{enumerate}

\textbf{Common checks and traps.} Sticky is OR of all discarded bits, not just one bit.

\subsection{Floating-Point Division Overview (L11)}
\textbf{What it is.} FP division computes $A/B$ by forming the reciprocal of $B$'s significand and multiplying.

\textbf{Step-by-step.}
\begin{enumerate}[leftmargin=*]
  \item Compute sign: $S = S_A \oplus S_B$.
  \item Exponent: $E_C = E_A - E_B + \text{bias}$.
  \item Compute reciprocal of significand $1/S_B$ using iteration.
  \item Multiply $S_A$ by $1/S_B$, normalize, and round.
\end{enumerate}

\textbf{Examples.}
\begin{enumerate}[leftmargin=*]
  \item Divide by power of two: only exponent changes.
  \item $A/B$ with $B$ in $[1,2)$ fits the reciprocal iteration range.
  \item $A/B$ with $B=1.0$ leaves significand unchanged.
\end{enumerate}

\textbf{Common checks and traps.} No remainder in FP division; treat as real division.

\subsection{Goldschmidt Division Method (L11)}
\textbf{What it is.} Computes reciprocal using multiplications and additions to approximate $1/S_B$.

\textbf{Step-by-step.}
\begin{enumerate}[leftmargin=*]
  \item Let $S_B = 1 + X$ where $0 < X < 1$.
  \item Transform to $1/(1-Y)$ with $Y = X/2$ and iterate $(1+Y)(1+Y^2)(1+Y^4)...$.
  \item Use FP multipliers to compute powers and accumulate.
  \item Multiply by $S_A$ and adjust exponent.
\end{enumerate}

\textbf{Examples.}
\begin{enumerate}[leftmargin=*]
  \item $S_B=1.11_2$ gives $X=0.11_2$ and $Y=0.001_2$.
  \item Four multiplications and five additions give a usable reciprocal for single precision.
  \item If $Y$ is small, higher powers like $Y^{32}$ become negligible.
\end{enumerate}

\textbf{Common checks and traps.} Ensure $S_B$ is normalized so $1 < S_B < 2$.

\subsection{Newton-Raphson Division Method (L11)}
\textbf{What it is.} Finds reciprocal using root finding on $f(x)=1/x-b$ with fast convergence.

\textbf{Step-by-step.}
\begin{enumerate}[leftmargin=*]
  \item Choose $x_0$ from a lookup table, with $0.5 < x_0 \le 1$.
  \item Iterate $x_{i+1} = x_i(2 - b x_i)$.
  \item Stop when precision matches target (e.g., 23 bits for single).
  \item Multiply by $S_A$ and normalize.
\end{enumerate}

\textbf{Examples.}
\begin{enumerate}[leftmargin=*]
  \item $b=1.5$, $x_0=0.75$ gives $x_1=0.75(2-1.125)=0.65625$.
  \item Error roughly squares each iteration; 5 iterations suffice for single precision.
  \item If $x_0$ is poor, more iterations are needed.
\end{enumerate}

\textbf{Common checks and traps.} Keep $b$ in $(1,2)$ by normalizing the input.

\subsection{FP Hardware: Multi-Cycle and Pipelined Units (L10, L11)}
\textbf{What it is.} FP adders and multipliers are more complex than integer units and usually take multiple cycles; they can be pipelined.

\textbf{Step-by-step.}
\begin{enumerate}[leftmargin=*]
  \item Compare exponents and align.
  \item Perform significand operation.
  \item Normalize and round.
  \item Detect overflow/underflow.
\end{enumerate}

\textbf{Examples.}
\begin{enumerate}[leftmargin=*]
  \item An FP adder may take several cycles; integer add is usually 1 cycle.
  \item Pipelined FP units allow throughput > 1 op per cycle.
  \item A slow FP unit sets the clock period if it is on the critical path.
\end{enumerate}

\textbf{Common checks and traps.} FP pipelines improve throughput, not single-op latency.

\subsection{MIPS FPU Architecture (L10)}
\textbf{What it is.} The FP unit is coprocessor 1 with its own 32 registers $f0$--$f31$.

\textbf{Step-by-step.}
\begin{enumerate}[leftmargin=*]
  \item Single precision uses one register.
  \item Double precision uses even-odd pairs ($f0/f1$, $f2/f3$, ...).
  \item FP ops only use FP registers.
\end{enumerate}

\textbf{Examples.}
\begin{enumerate}[leftmargin=*]
  \item Use $f2$ for a single precision value and $f4/f5$ for a double.
  \item $add.s f0, f1, f6$ is a single-precision add.
  \item $mul.d f4, f4, f6$ uses double-precision pairs.
\end{enumerate}

\textbf{Common checks and traps.} Double precision must start at even register numbers.

\subsection{MIPS FP Instructions and Data Movement (L10)}
\textbf{What it is.} Separate FP arithmetic, load/store, move, convert, and compare instructions.

\textbf{Step-by-step.}
\begin{enumerate}[leftmargin=*]
  \item Use \texttt{add.s/sub.s/mul.s/div.s} for single.
  \item Use \texttt{add.d/sub.d/mul.d/div.d} for double.
  \item Use \texttt{lwc1/swc1} and \texttt{ldc1/sdc1} for loads/stores.
  \item Use \texttt{mfc1/mtc1} to move between GPR and FPR.
  \item Use \texttt{c.eq.s/c.lt.s/c.le.s} and \texttt{bc1t/bc1f} for FP branches.
\end{enumerate}

\textbf{Examples.}
\begin{enumerate}[leftmargin=*]
  \item Load, add, store: \texttt{lwc1 f2, 0(t0); lwc1 f4, 4(t0); add.s f6, f2, f4; swc1 f6, 8(t0)}.
  \item Convert int to float: \texttt{mtc1 t0, f0; cvt.s.w f2, f0}.
  \item Compare and branch: \texttt{c.lt.s f2, f4; bc1t label}.
\end{enumerate}

\textbf{Common checks and traps.} FP branches use the FP condition flag, not integer flags.

\subsection{FP Syscalls and Example (L10)}
\textbf{What it is.} MARS syscalls include read/print float and double, and FP arguments use $f12$.

\textbf{Step-by-step.}
\begin{enumerate}[leftmargin=*]
  \item Set $v0=2$ to print float, $v0=3$ for double.
  \item Place value in $f12$ and call \texttt{syscall}.
  \item For reading, use $v0=6$ (float) or $v0=7$ (double).
\end{enumerate}

\textbf{Examples.}
\begin{enumerate}[leftmargin=*]
  \item Print float: \texttt{li v0,2; mov.s f12,f0; syscall}.
  \item Read float: \texttt{li v0,6; syscall} (result in $f0$).
  \item Fahrenheit to Celsius: $C = (5.0/9.0)(F-32.0)$ using \texttt{sub.s, div.s, mul.s}.
\end{enumerate}

\textbf{Common checks and traps.} $f12$ is the argument register for FP prints.

\subsection{Integer Multiplication Complexity (L10)}
\textbf{What it is.} Classical and faster multiplication algorithms have different asymptotic costs.

\textbf{Step-by-step.}
\begin{enumerate}[leftmargin=*]
  \item School method: $\Theta(n^2)$ bit operations.
  \item Karatsuba: reduce 4 multiplies to 3, $O(n^{\log_2 3}) \approx O(n^{1.585})$.
  \item Schonhage--Strassen: $O(n \log n \log \log n)$.
\end{enumerate}

\textbf{Examples.}
\begin{enumerate}[leftmargin=*]
  \item For $n=1024$, Karatsuba beats $n^2$ for large enough $n$.
  \item Karatsuba splits 16-bit numbers into two 8-bit halves.
  \item As $n$ grows, $n \log n \log \log n$ becomes far smaller than $n^2$.
\end{enumerate}

\textbf{Common checks and traps.} Asymptotic improvement appears only for large $n$.

\section{Integer Arithmetic and ALU (Lectures 11--12)}

\subsection{Unsigned Shift-Add Multiplier (L12)}
\textbf{What it is.} Hardware multiplication by repeated add and shift using a 64-bit product register.

\textbf{Step-by-step.}
\begin{enumerate}[leftmargin=*]
  \item Initialize product with multiplier in the low half, zeros in the high half.
  \item If product LSB is 1, add multiplicand to the high half.
  \item Shift the entire product right by 1.
  \item Repeat for word width (e.g., 32 iterations).
\end{enumerate}

\textbf{Examples.}
\begin{enumerate}[leftmargin=*]
  \item $0010_2 \times 0110_2 = 1100_2$ via four iterations.
  \item $1000_2 \times 1001_2$ adds only on cycles where LSB is 1.
  \item Multiply by 0 yields zero with no effective adds.
\end{enumerate}

\textbf{Common checks and traps.} The add uses the high half of the product, not the low.

\begin{figure}[h]
\centering
\fbox{\parbox{0.9\linewidth}{Diagram placeholder: Shift-add multiplier datapath (L12)}}
\end{figure}

\subsection{Signed Multiplication and Sign Handling (L12)}
\textbf{What it is.} Signed product sign is XOR of operand signs; multiply magnitudes, then fix sign.

\textbf{Step-by-step.}
\begin{enumerate}[leftmargin=*]
  \item Compute sign of result: $S = S_A \oplus S_B$.
  \item Convert operands to positive magnitudes.
  \item Multiply as unsigned.
  \item If $S=1$, negate the product.
\end{enumerate}

\textbf{Examples.}
\begin{enumerate}[leftmargin=*]
  \item $(-3) \times (+2) = -6$.
  \item $(-3) \times (-2) = +6$.
  \item $(-8) \times 0 = 0$ (sign irrelevant).
\end{enumerate}

\textbf{Common checks and traps.} Ensure operands are converted using two's complement rules.

\subsection{Booth's Encoding and Algorithm (L12)}
\textbf{What it is.} Booth's algorithm reduces runs of 1s in the multiplier to fewer add/sub operations.

\textbf{Step-by-step.}
\begin{enumerate}[leftmargin=*]
  \item Look at current bit and the bit to the right (previous bit).
  \item 10: subtract multiplicand; 01: add multiplicand; 00 or 11: no op.
  \item Arithmetic shift right after each step.
\end{enumerate}

\textbf{Examples.}
\begin{enumerate}[leftmargin=*]
  \item Multiplier 01111 gives $16x - x$ (one subtract, shifts).
  \item Multiplier 1101 triggers a subtract at 10 and add at 01.
  \item A run of 1s uses two operations, not one per bit.
\end{enumerate}

\textbf{Common checks and traps.} Use arithmetic shift for signed values.

\begin{figure}[h]
\centering
\fbox{\parbox{0.9\linewidth}{Diagram placeholder: Booth encoding table and example (L12)}}
\end{figure}

\subsection{Full Adder, Ripple-Carry Adder, and 2's Complement Subtraction (L11, L12)}
\textbf{What it is.} A full adder computes sum and carry; ripple-carry chains adders; subtraction uses two's complement.

\textbf{Step-by-step.}
\begin{enumerate}[leftmargin=*]
  \item Full adder: $\text{Sum}=A\oplus B\oplus C_{in}$, $C_{out}=AB+BC_{in}+AC_{in}$.
  \item Ripple-carry: chain $n$ full adders; carry propagates left.
  \item Subtraction: invert $B$ and set $C_{in}=1$ to add $A + (\sim B + 1)$.
\end{enumerate}

\textbf{Examples.}
\begin{enumerate}[leftmargin=*]
  \item 1-bit: $A=1, B=1, C_{in}=0$ gives Sum=0, Carry=1.
  \item 4-bit RCA delay is about 4 full-adder delays.
  \item $7-3$ in 4-bit two's complement: $0111 + 1101 = 0100$.
\end{enumerate}

\textbf{Common checks and traps.} RCA delay is $O(n)$, which sets the clock.

\subsection{Carry Lookahead Adder (CLA) (L11)}
\textbf{What it is.} CLA uses generate/propagate to compute carries in parallel.

\textbf{Step-by-step.}
\begin{enumerate}[leftmargin=*]
  \item Generate: $g_i = a_i b_i$.
  \item Propagate: $p_i = a_i + b_i$.
  \item Carry: $c_{i+1} = g_i + p_i c_i$.
  \item Expand for 4-bit group to compute $c_1..c_4$ directly.
\end{enumerate}

\textbf{Examples.}
\begin{enumerate}[leftmargin=*]
  \item Compute $c_1 = g_0 + p_0 c_0$ for given bits.
  \item For a 4-bit group, $c_4$ uses $g_3, p_3, g_2, ...$.
  \item CLA reduces delay compared to RCA at the cost of extra logic.
\end{enumerate}

\textbf{Common checks and traps.} Use OR for propagate in these slides, not XOR.

\subsection{Hybrid CLA/RCA Adders (L11)}
\textbf{What it is.} Combine fast CLA within groups and ripple between groups to balance cost and speed.

\textbf{Step-by-step.}
\begin{enumerate}[leftmargin=*]
  \item Split the adder into groups (e.g., 4-bit groups).
  \item Use CLA inside each group to compute local carries.
  \item Ripple the group carries between groups.
\end{enumerate}

\textbf{Examples.}
\begin{enumerate}[leftmargin=*]
  \item 12-bit adder as three 4-bit CLA groups.
  \item Larger groups reduce ripple levels but increase logic cost.
  \item Compare 12-bit pure RCA vs hybrid CLA+RCA delay.
\end{enumerate}

\textbf{Common checks and traps.} Full CLA for large n has impractical cost.

\subsection{32-bit MIPS ALU Design (L12)}
\textbf{What it is.} A 32-bit ALU built from 1-bit slices supports AND, OR, ADD, SUB, and SLT.

\textbf{Step-by-step.}
\begin{enumerate}[leftmargin=*]
  \item Use control lines (Binvert, Op1, Op0) to select operation.
  \item For SUB, invert B and set carry-in to 1.
  \item For SLT, use the sign of $A-B$ and set bit 0.
\end{enumerate}

\textbf{Examples.}
\begin{enumerate}[leftmargin=*]
  \item AND: Binvert=0, Op=00.
  \item ADD: Binvert=0, Op=10.
  \item SUB: Binvert=1, Op=10.
\end{enumerate}

\textbf{Common checks and traps.} SLT is based on signed comparison of $A-B$.

\subsection{Overflow Detection and SLT Support (L12)}
\textbf{What it is.} Signed overflow in add/sub occurs when carry into MSB differs from carry out.

\textbf{Step-by-step.}
\begin{enumerate}[leftmargin=*]
  \item Compute carries into and out of bit 31.
  \item Overflow = carry31 XOR carry32.
  \item SLT uses the sign of the subtraction result with overflow adjustment.
\end{enumerate}

\textbf{Examples.}
\begin{enumerate}[leftmargin=*]
  \item $0x7FFFFFFF + 1$ overflows to negative.
  \item $0x80000000 - 1$ overflows to positive.
  \item SLT sets result to 1 if $A < B$ else 0.
\end{enumerate}

\textbf{Common checks and traps.} Overflow rules differ for signed vs unsigned adds.

\subsection{MIPS Instruction Execution Cycle (L12)}
\textbf{What it is.} Classic 5-stage flow: IF, ID/RR, EX, MEM, WB.

\textbf{Step-by-step.}
\begin{enumerate}[leftmargin=*]
  \item IF: fetch instruction and compute PC+4.
  \item ID/RR: decode and read registers.
  \item EX: ALU operation or address calculation.
  \item MEM: data memory access for loads/stores.
  \item WB: write result back to register file.
\end{enumerate}

\textbf{Examples.}
\begin{enumerate}[leftmargin=*]
  \item R-type: uses IF, ID, EX, WB (MEM is idle).
  \item lw: uses all 5 stages.
  \item beq: uses IF, ID, EX (MEM/WB idle in single-cycle view).
\end{enumerate}

\textbf{Common checks and traps.} Different instruction types still share the same stage order.

\subsection{Combinational vs Sequential Elements and Register File (L12)}
\textbf{What it is.} Combinational logic has no memory (ALU, mux). Sequential elements store state (registers).

\textbf{Step-by-step.}
\begin{enumerate}[leftmargin=*]
  \item Register file reads two registers and writes one per cycle.
  \item Writes occur on the clock edge; reads are combinational.
  \item Control signals gate the write operation.
\end{enumerate}

\textbf{Examples.}
\begin{enumerate}[leftmargin=*]
  \item A register holds data across cycles using D flip-flops.
  \item Read rs and rt while writing rd in the same cycle.
  \item Write enable 0 leaves register unchanged.
\end{enumerate}

\textbf{Common checks and traps.} Writes are edge-triggered; do not assume immediate updates mid-cycle.

\section{Processor Design: Single-Cycle MIPS (Lectures 13--14)}

\subsection{ISA-Driven Datapath Design and Building Blocks (L13)}
\textbf{What it is.} Start from ISA requirements and map to hardware blocks: ALU, registers, muxes, memory, adders.

\textbf{Step-by-step.}
\begin{enumerate}[leftmargin=*]
  \item Identify instruction types (R, load/store, branch, jump).
  \item List required operations per instruction (read reg, ALU, memory, write reg).
  \item Choose hardware blocks to cover all operations.
  \item Connect blocks with multiplexers for shared paths.
\end{enumerate}

\textbf{Examples.}
\begin{enumerate}[leftmargin=*]
  \item R-type needs register read, ALU, register write.
  \item Load needs ALU for address and memory read.
  \item Branch needs register compare and PC update.
\end{enumerate}

\textbf{Common checks and traps.} Avoid assuming a block can do two functions in one cycle.

\begin{figure}[h]
\centering
\fbox{\parbox{0.9\linewidth}{Diagram placeholder: Single-cycle datapath overview with control signals (L13--L14)}}
\end{figure}

\subsection{R-Type Datapath (L13)}
\textbf{What it is.} Reg-reg instructions read two registers, compute in ALU, and write to rd.

\textbf{Step-by-step.}
\begin{enumerate}[leftmargin=*]
  \item Read rs and rt from register file.
  \item Select ALU operation based on funct.
  \item Write result to rd.
\end{enumerate}

\textbf{Examples.}
\begin{enumerate}[leftmargin=*]
  \item add rd, rs, rt: RegDst=1, ALUSrc=0, RegWrite=1.
  \item and rd, rs, rt: ALUOp selects AND.
  \item slt rd, rs, rt: ALU computes set-less-than.
\end{enumerate}

\textbf{Common checks and traps.} rd is destination for R-type, not rt.

\subsection{Load Datapath (L13)}
\textbf{What it is.} lw computes address in ALU, reads memory, writes to rt.

\textbf{Step-by-step.}
\begin{enumerate}[leftmargin=*]
  \item Read rs and sign-extend immediate.
  \item ALU adds rs + immediate for address.
  \item Read memory and write data to rt.
\end{enumerate}

\textbf{Examples.}
\begin{enumerate}[leftmargin=*]
  \item lw rt, 16(rs): ALUSrc=1, MemRead=1, MemtoReg=1.
  \item Base+offset address uses signed immediate.
  \item lw uses RegDst=0 because rt is destination.
\end{enumerate}

\textbf{Common checks and traps.} Immediate must be sign-extended for address calculation.

\subsection{Store Datapath (L13)}
\textbf{What it is.} sw computes address in ALU and writes register value to memory.

\textbf{Step-by-step.}
\begin{enumerate}[leftmargin=*]
  \item Read rs (base) and rt (data).
  \item ALU adds rs + immediate for address.
  \item Write rt data to memory.
\end{enumerate}

\textbf{Examples.}
\begin{enumerate}[leftmargin=*]
  \item sw rt, 0(rs): MemWrite=1, RegWrite=0.
  \item Store uses ALUSrc=1 to add offset.
  \item sw does not write back to register file.
\end{enumerate}

\textbf{Common checks and traps.} sw uses rt as source, not destination.

\subsection{Branch Datapath (L13)}
\textbf{What it is.} beq compares rs and rt and conditionally updates PC.

\textbf{Step-by-step.}
\begin{enumerate}[leftmargin=*]
  \item Read rs and rt.
  \item ALU subtracts to set Zero flag.
  \item Sign-extend and shift immediate left 2.
  \item Add to PC+4 for target address.
\end{enumerate}

\textbf{Examples.}
\begin{enumerate}[leftmargin=*]
  \item beq rs, rt, imm: Branch=1, ALUOp=sub.
  \item If Zero=1, PCsrc selects branch target.
  \item If Zero=0, PC continues with PC+4.
\end{enumerate}

\textbf{Common checks and traps.} Branch offset is word-aligned, so shift left by 2.

\subsection{Combining Datapaths with Multiplexers (L13)}
\textbf{What it is.} Different instructions reuse hardware with muxes to select inputs and outputs.

\textbf{Step-by-step.}
\begin{enumerate}[leftmargin=*]
  \item Use ALUSrc mux to choose register vs immediate.
  \item Use RegDst mux to choose rd vs rt.
  \item Use MemtoReg mux to choose ALU vs memory data.
  \item Use PCsrc mux for PC+4 vs branch target.
\end{enumerate}

\textbf{Examples.}
\begin{enumerate}[leftmargin=*]
  \item lw uses ALUSrc=1 and MemtoReg=1.
  \item R-type uses RegDst=1 and MemtoReg=0.
  \item beq uses PCsrc=1 only if Zero=1.
\end{enumerate}

\textbf{Common checks and traps.} Do not wire two sources together without a mux.

\subsection{Control Signals for Single-Cycle MIPS (L14)}
\textbf{What it is.} Control signals drive muxes and memory operations based on opcode.

\textbf{Step-by-step.}
\begin{enumerate}[leftmargin=*]
  \item Decode opcode to generate RegDst, ALUSrc, MemRead, MemWrite, MemtoReg, RegWrite, Branch, ALUOp.
  \item Set control values for each instruction type.
  \item Feed ALUOp to ALU control to select function.
\end{enumerate}

\textbf{Examples.}
\begin{enumerate}[leftmargin=*]
  \item R-type: RegDst=1, ALUSrc=0, RegWrite=1, MemRead=0, MemWrite=0.
  \item lw: RegDst=0, ALUSrc=1, MemRead=1, MemtoReg=1, RegWrite=1.
  \item sw: ALUSrc=1, MemWrite=1, RegWrite=0.
\end{enumerate}

\textbf{Common checks and traps.} Branch control also depends on Zero output.

\subsection{ALU Control Logic (L14)}
\textbf{What it is.} Combines ALUOp and funct to select ALU operation.

\textbf{Step-by-step.}
\begin{enumerate}[leftmargin=*]
  \item ALUOp=00 for load/store (add).
  \item ALUOp=01 for beq (subtract).
  \item ALUOp=10 for R-type (use funct).
  \item Map funct to AND/OR/ADD/SUB/SLT codes.
\end{enumerate}

\textbf{Examples.}
\begin{enumerate}[leftmargin=*]
  \item funct 100000 -> add.
  \item funct 100010 -> subtract.
  \item funct 101010 -> slt.
\end{enumerate}

\textbf{Common checks and traps.} ALUOp only distinguishes classes; funct picks final op.

\subsection{Main Control Unit Design (L14)}
\textbf{What it is.} A combinational controller that outputs control signals from opcode bits.

\textbf{Step-by-step.}
\begin{enumerate}[leftmargin=*]
  \item Use opcode bits as inputs to logic.
  \item Produce control signals for datapath.
  \item Verify each instruction's control vector.
\end{enumerate}

\textbf{Examples.}
\begin{enumerate}[leftmargin=*]
  \item For opcode 0 (R-type), RegDst=1 and ALUOp=10.
  \item For opcode 35 (lw), MemRead=1 and MemtoReg=1.
  \item For opcode 4 (beq), Branch=1 and ALUOp=01.
\end{enumerate}

\textbf{Common checks and traps.} Single-cycle control is purely combinational.

\subsection{Jump Instruction Handling (L14)}
\textbf{What it is.} j uses a 26-bit target combined with upper PC bits.

\textbf{Step-by-step.}
\begin{enumerate}[leftmargin=*]
  \item Take upper 4 bits of PC+4.
  \item Concatenate 26-bit target and 2 zeros (word aligned).
  \item Select this as the next PC when jump opcode detected.
\end{enumerate}

\textbf{Examples.}
\begin{enumerate}[leftmargin=*]
  \item Jump to 0x00401234 uses top 4 bits of PC+4.
  \item j does not read registers.
  \item j has a dedicated PC mux select line.
\end{enumerate}

\textbf{Common checks and traps.} Jump address is pseudo-absolute within current 256MB region.

\subsection{Critical Path and Performance Equation (L14)}
\textbf{What it is.} Single-cycle CCT is set by the longest instruction path, typically load.

\textbf{Step-by-step.}
\begin{enumerate}[leftmargin=*]
  \item Find the longest datapath: IMem -> RegFile -> ALU -> DMem -> RegFile.
  \item Set CCT to cover that path.
  \item CPU time: $\text{IC} \times \text{CPI} \times \text{CCT}$.
\end{enumerate}

\textbf{Examples.}
\begin{enumerate}[leftmargin=*]
  \item Load has the longest path, so it sets clock period.
  \item Fast R-type instructions still wait for long CCT.
  \item Pipelining lowers CCT at the cost of higher CPI.
\end{enumerate}

\textbf{Common checks and traps.} Single-cycle design violates "make the common case fast".

\subsection{Logic Faults and Stuck-at Errors (L14)}
\textbf{What it is.} A signal line stuck at 0 or 1 can cause certain instructions to fail.

\textbf{Step-by-step.}
\begin{enumerate}[leftmargin=*]
  \item Identify which control or data line is stuck.
  \item Determine which instructions depend on that line.
  \item Evaluate functional impact (safe vs catastrophic).
\end{enumerate}

\textbf{Examples.}
\begin{enumerate}[leftmargin=*]
  \item If ALUSrc stuck at 0, lw/sw will fail (need immediate).
  \item If RegWrite stuck at 0, all reg writes fail.
  \item If Branch control stuck at 1, PC updates incorrectly.
\end{enumerate}

\textbf{Common checks and traps.} Some faults break only a subset of instructions.

\section{Pipelining and Hazards (Lectures 15--18)}

\subsection{Pipelining Basics: Throughput vs Latency (L15)}
\textbf{What it is.} Pipelining overlaps instruction stages to increase throughput without reducing single-instruction latency.

\textbf{Step-by-step.}
\begin{enumerate}[leftmargin=*]
  \item Split execution into stages (IF, ID, EX, MEM, WB).
  \item Start a new instruction each cycle.
  \item After pipeline fills, complete roughly one instruction per cycle.
\end{enumerate}

\textbf{Examples.}
\begin{enumerate}[leftmargin=*]
  \item Laundry analogy: 4 loads complete faster with pipeline than sequential.
  \item A 5-stage pipeline ideally gives near 5x throughput for long programs.
  \item Latency of one instruction still spans 5 stages.
\end{enumerate}

\textbf{Common checks and traps.} Pipelining boosts throughput, not single-op speed.

\subsection{MIPS 5-Stage Pipeline and Diagram (L15)}
\textbf{What it is.} MIPS pipeline stages: IF, ID/RR, EX, MEM, WB.

\textbf{Step-by-step.}
\begin{enumerate}[leftmargin=*]
  \item IF: fetch instruction and PC+4.
  \item ID/RR: decode and read registers.
  \item EX: ALU operation or address calculation.
  \item MEM: access data memory.
  \item WB: write result back.
\end{enumerate}

\textbf{Examples.}
\begin{enumerate}[leftmargin=*]
  \item Pipeline diagram shows 5 instructions active at cycle 5.
  \item R-type uses MEM as NOP stage for uniformity.
  \item Load uses all stages; store uses MEM but no WB.
\end{enumerate}

\textbf{Common checks and traps.} All instructions occupy the same number of stages (with NOPs).

\begin{figure}[h]
\centering
\fbox{\parbox{0.9\linewidth}{Diagram placeholder: 5-stage pipeline timeline (L15)}}
\end{figure}

\subsection{Pipeline Terminology: Fill, Full, Drain, Balance (L15)}
\textbf{What it is.} Pipeline fills, becomes full, then drains. Speedup is limited by the slowest stage.

\textbf{Step-by-step.}
\begin{enumerate}[leftmargin=*]
  \item Fill: initial cycles before all stages are active.
  \item Full: steady state, one completion per cycle.
  \item Drain: final cycles after last instruction enters.
  \item Balance stages to minimize cycle time.
\end{enumerate}

\textbf{Examples.}
\begin{enumerate}[leftmargin=*]
  \item 5-stage pipeline becomes full at cycle 5.
  \item If EX is longest stage, CCT equals EX time.
  \item Unbalanced stages reduce speedup below pipeline depth.
\end{enumerate}

\textbf{Common checks and traps.} Speedup $< k$ even without hazards.

\subsection{Pipeline Registers and Overhead (L15, L16)}
\textbf{What it is.} Latches between stages hold intermediate values; they add overhead.

\textbf{Step-by-step.}
\begin{enumerate}[leftmargin=*]
  \item Insert registers between each pair of stages.
  \item Propagate instruction bits and data each cycle.
  \item Account for latch delay in CCT.
\end{enumerate}

\textbf{Examples.}
\begin{enumerate}[leftmargin=*]
  \item IF/ID holds instruction and PC+4.
  \item EX/MEM holds ALU result and control bits.
  \item Latch overhead reduces ideal speedup.
\end{enumerate}

\textbf{Common checks and traps.} Pipeline registers must carry control signals too.

\subsection{Structural Hazards (L16)}
\textbf{What it is.} Two stages need the same resource in the same cycle.

\textbf{Step-by-step.}
\begin{enumerate}[leftmargin=*]
  \item Identify shared resources (e.g., single memory).
  \item If conflict, stall one instruction (bubble).
  \item Prefer separate instruction and data memories.
\end{enumerate}

\textbf{Examples.}
\begin{enumerate}[leftmargin=*]
  \item IF and MEM both need memory in the same cycle.
  \item Separate I-cache and D-cache removes the conflict.
  \item Multi-ported memory is another solution.
\end{enumerate}

\textbf{Common checks and traps.} Structural hazards are solved by hardware duplication or stalls.

\subsection{Data Hazards: RAW, WAR, WAW (L16, L17)}
\textbf{What it is.} An instruction depends on data from a previous instruction that is not yet written.

\textbf{Step-by-step.}
\begin{enumerate}[leftmargin=*]
  \item RAW: read after write (most common).
  \item WAW: write after write (ordering issue).
  \item WAR: write after read (possible in out-of-order).
\end{enumerate}

\textbf{Examples.}
\begin{enumerate}[leftmargin=*]
  \item RAW: \texttt{add \$t1, ...} then sub uses \texttt{\$t1}.
  \item WAW: \texttt{lw} writes \texttt{\$t1}, then \texttt{sub} writes \texttt{\$t1} earlier.
  \item WAR: out-of-order write overtakes read.
\end{enumerate}

\textbf{Common checks and traps.} Classic in-order MIPS has RAW hazards; WAR/WAW are rare.

\subsection{Half-Cycle Register Write/Read Fix (L16)}
\textbf{What it is.} Write in first half-cycle, read in second half-cycle to reduce RAW stalls.

\textbf{Step-by-step.}
\begin{enumerate}[leftmargin=*]
  \item Write back at beginning of cycle.
  \item Allow ID/RR to read updated value in same cycle.
  \item This removes some, not all, RAW hazards.
\end{enumerate}

\textbf{Examples.}
\begin{enumerate}[leftmargin=*]
  \item \texttt{add} writes \texttt{\$t1} in WB half-cycle; next instruction reads in ID.
  \item Removes one stall in some sequences.
  \item Still fails for load-use when value not ready.
\end{enumerate}

\textbf{Common checks and traps.} Load data arrives after MEM, so half-cycle does not fix load-use.

\subsection{Load-Use Hazards and Bubbles (L16, L17)}
\textbf{What it is.} A load followed immediately by a consumer needs a stall cycle.

\textbf{Step-by-step.}
\begin{enumerate}[leftmargin=*]
  \item Detect that the next instruction uses the load destination.
  \item Insert a bubble (NOP) to delay the consumer.
  \item Continue once data is available.
\end{enumerate}

\textbf{Examples.}
\begin{enumerate}[leftmargin=*]
  \item \texttt{lw \$t1, 0(\$t2)} then \texttt{add \$t3, \$t1, \$t4} requires 1 stall.
  \item With two consumers, multiple stalls may be needed.
  \item Inserting a useful independent instruction can hide the stall.
\end{enumerate}

\textbf{Common checks and traps.} Forwarding cannot bring data backward in time.

\subsection{Forwarding (Bypassing) (L16, L17)}
\textbf{What it is.} Use ALU or MEM results directly as ALU inputs for later instructions.

\textbf{Step-by-step.}
\begin{enumerate}[leftmargin=*]
  \item Add EX->EX and MEM->EX bypass paths.
  \item Use forwarding control to select newest value.
  \item Eliminate most RAW hazards.
\end{enumerate}

\textbf{Examples.}
\begin{enumerate}[leftmargin=*]
  \item add then sub uses forwarded EX result.
  \item MEM->EX forwards load data once in MEM/WB.
  \item With forwarding, many sequences drop from 10 to 8 cycles.
\end{enumerate}

\textbf{Common checks and traps.} Load-use still needs 1 stall even with forwarding.

\begin{figure}[h]
\centering
\fbox{\parbox{0.9\linewidth}{Diagram placeholder: Forwarding paths and hazard detection (L17)}}
\end{figure}

\subsection{Software Interlocking and Scheduling (L17)}
\textbf{What it is.} Compiler reorders instructions or inserts NOPs to avoid hazards.

\textbf{Step-by-step.}
\begin{enumerate}[leftmargin=*]
  \item Identify dependent instructions.
  \item Move independent instructions into delay slots.
  \item Insert NOPs when no safe instruction exists.
\end{enumerate}

\textbf{Examples.}
\begin{enumerate}[leftmargin=*]
  \item Swap a load with a later independent add.
  \item Insert a NOP after load when no independent work exists.
  \item Reorder loop body to hide load latency.
\end{enumerate}

\textbf{Common checks and traps.} Reordering must preserve program semantics.

\subsection{Control Hazards and Branch Penalty (L17, L18)}
\textbf{What it is.} The next PC is unknown until branch outcome is known.

\textbf{Step-by-step.}
\begin{enumerate}[leftmargin=*]
  \item In classic MIPS, branch target is resolved in MEM stage.
  \item This causes a 3-cycle branch penalty.
  \item Stall or flush instructions until outcome is known.
\end{enumerate}

\textbf{Examples.}
\begin{enumerate}[leftmargin=*]
  \item beq followed by lw: lw may be squashed if branch taken.
  \item 3 stalls per branch if no prediction is used.
  \item Moving compare earlier reduces penalty to 1.
\end{enumerate}

\textbf{Common checks and traps.} PC update point determines penalty length.

\begin{figure}[h]
\centering
\fbox{\parbox{0.9\linewidth}{Diagram placeholder: Branch hazard pipeline timing (L18)}}
\end{figure}

\subsection{Early Branch Resolution (L18)}
\textbf{What it is.} Move compare and target computation to ID/RR to reduce stalls.

\textbf{Step-by-step.}
\begin{enumerate}[leftmargin=*]
  \item Add a zero-checker in ID.
  \item Compute branch target in ID.
  \item Update PC earlier if branch is taken.
\end{enumerate}

\textbf{Examples.}
\begin{enumerate}[leftmargin=*]
  \item Reduces branch penalty from 3 to 1 cycle.
  \item Can create new hazards if branch depends on prior load.
  \item Needs extra ALU hardware for ID stage.
\end{enumerate}

\textbf{Common checks and traps.} Early branch logic can conflict with earlier instruction writes.

\subsection{Branch Prediction (L18)}
\textbf{What it is.} Predict branch outcome to avoid stalls; correct prediction avoids penalty.

\textbf{Step-by-step.}
\begin{enumerate}[leftmargin=*]
  \item Predict-not-taken: fetch PC+4 by default.
  \item Static prediction: backward taken, forward not-taken.
  \item Dynamic prediction: keep history per branch.
\end{enumerate}

\textbf{Examples.}
\begin{enumerate}[leftmargin=*]
  \item 47\% of MIPS branches are not taken on average.
  \item A loop branch is usually taken; static predictor exploits this.
  \item Dynamic prediction stores last outcome in a history table.
\end{enumerate}

\textbf{Common checks and traps.} Wrong prediction requires pipeline flush.

\subsection{1-bit vs 2-bit Predictors (L18)}
\textbf{What it is.} 2-bit predictors reduce flip-flopping on loop exits.

\textbf{Step-by-step.}
\begin{enumerate}[leftmargin=*]
  \item 1-bit predictor flips on every misprediction.
  \item 2-bit predictor changes prediction only after two misses.
  \item Improves accuracy for loop branches.
\end{enumerate}

\textbf{Examples.}
\begin{enumerate}[leftmargin=*]
  \item Inner loop exit causes two mispredictions with 1-bit predictor.
  \item 2-bit predictor stays in "taken" state through one miss.
  \item For random branches, 2-bit may not help much.
\end{enumerate}

\textbf{Common checks and traps.} Prediction state must be updated on actual outcome.

\subsection{Delay Slots (L18)}
\textbf{What it is.} Compiler places a useful instruction after a branch to fill the delay slot.

\textbf{Step-by-step.}
\begin{enumerate}[leftmargin=*]
  \item Identify the branch delay slot.
  \item Move a safe instruction into that slot.
  \item If none exists, insert NOP.
\end{enumerate}

\textbf{Examples.}
\begin{enumerate}[leftmargin=*]
  \item Fill with an instruction from fall-through path.
  \item Fill with an instruction from taken path if safe.
  \item Use NOP if reordering breaks correctness.
\end{enumerate}

\textbf{Common checks and traps.} Delay-slot instruction always executes, taken or not.

\subsection{Pipeline Performance with Stalls (L18)}
\textbf{What it is.} Pipeline speedup drops with stall cycles.

\textbf{Step-by-step.}
\begin{enumerate}[leftmargin=*]
  \item Ideal CPI for pipeline is 1.
  \item Actual CPI = $1 + \text{avg stall cycles per instruction}$.
  \item Speedup $= \text{pipe depth} / \text{CPI}$.
\end{enumerate}

\textbf{Examples.}
\begin{enumerate}[leftmargin=*]
  \item With branch frequency 0.2 and penalty 1: CPI = 1 + 0.2.
  \item With load-use stalls 0.1 per instruction: CPI = 1.1.
  \item Speedup of 5-stage pipeline with CPI 1.25 is 4.0.
\end{enumerate}

\textbf{Common checks and traps.} CPI uses average stalls, not worst-case.

\section{Memory Hierarchy, Caches, and Virtual Memory (Lectures 18--20)}

\subsection{Memory Hierarchy Levels and Tradeoffs (L18, L20)}
\textbf{What it is.} Small, fast memories are costly; large memories are slow. Hierarchy bridges the gap.

\textbf{Step-by-step.}
\begin{enumerate}[leftmargin=*]
  \item Registers and L1 cache are fastest and smallest.
  \item Lower levels (L2, L3, DRAM, disk) are larger and slower.
  \item Data moves up and down the hierarchy in blocks.
\end{enumerate}

\textbf{Examples.}
\begin{enumerate}[leftmargin=*]
  \item L1 cache is tens of KB; DRAM is GBs.
  \item Disk is orders of magnitude slower than DRAM.
  \item Tertiary storage (tape/cloud) is largest and slowest.
\end{enumerate}

\textbf{Common checks and traps.} Faster memory almost always costs more per bit.

\begin{figure}[h]
\centering
\fbox{\parbox{0.9\linewidth}{Diagram placeholder: Memory hierarchy pyramid (L18)}}
\end{figure}

\subsection{Principle of Locality (L18)}
\textbf{What it is.} Programs reuse data nearby in time (temporal) and space (spatial).

\textbf{Step-by-step.}
\begin{enumerate}[leftmargin=*]
  \item Temporal locality: recent data likely used again.
  \item Spatial locality: nearby addresses likely accessed soon.
  \item Caches exploit both by storing recent blocks.
\end{enumerate}

\textbf{Examples.}
\begin{enumerate}[leftmargin=*]
  \item Loop counter reused each iteration (temporal).
  \item Sequential array access (spatial).
  \item Instruction stream has strong spatial locality.
\end{enumerate}

\textbf{Common checks and traps.} Locality is empirical, not guaranteed.

\subsection{Cache Basics: Hit, Miss, Block Size, I-Cache and D-Cache (L18)}
\textbf{What it is.} Cache stores blocks from lower memory. Hits are fast; misses incur penalty.

\textbf{Step-by-step.}
\begin{enumerate}[leftmargin=*]
  \item CPU access checks cache first.
  \item On hit: serve from cache.
  \item On miss: fetch block from lower level.
  \item I-cache holds instructions; D-cache holds data.
\end{enumerate}

\textbf{Examples.}
\begin{enumerate}[leftmargin=*]
  \item Hit time 1 cycle, miss penalty 100 cycles.
  \item A 16-word block fetches adjacent words.
  \item Separate I-cache and D-cache avoid structural hazards.
\end{enumerate}

\textbf{Common checks and traps.} Block size is in words/bytes, not bits.

\subsection{Direct-Mapped Cache Mapping, Tags, and Valid Bit (L18, L19)}
\textbf{What it is.} Each memory block maps to exactly one cache line, identified by index and tag.

\textbf{Step-by-step.}
\begin{enumerate}[leftmargin=*]
  \item Split address into tag, index, and offset.
  \item Use index to select cache line.
  \item Compare tag and valid bit to decide hit/miss.
\end{enumerate}

\textbf{Examples.}
\begin{enumerate}[leftmargin=*]
  \item 5-bit address, 8-line cache: 3-bit index.
  \item Valid bit starts at 0 to avoid garbage hits.
  \item Tag mismatch causes conflict miss.
\end{enumerate}

\textbf{Common checks and traps.} Index is typically low-order bits above offset.

\begin{figure}[h]
\centering
\fbox{\parbox{0.9\linewidth}{Diagram placeholder: Direct-mapped address split and mapping (L18, L19)}}
\end{figure}

\subsection{Direct-Mapped Cache Example (L18)}
\textbf{What it is.} Step-by-step trace of hits and misses for a direct-mapped cache.

\textbf{Step-by-step.}
\begin{enumerate}[leftmargin=*]
  \item Start with all valid bits 0.
  \item Process each address; compute index and tag.
  \item Update cache line on misses.
\end{enumerate}

\textbf{Examples.}
\begin{enumerate}[leftmargin=*]
  \item Address sequence 22, 26, 22, 26, 16, 3, 16, 18.
  \item First accesses are misses; repeats may hit if not evicted.
  \item Conflicts occur when two blocks map to same index.
\end{enumerate}

\textbf{Common checks and traps.} A hit requires both tag match and valid=1.

\subsection{Cache Miss Types (L19, L20)}
\textbf{What it is.} Misses are classified as compulsory, capacity, or conflict.

\textbf{Step-by-step.}
\begin{enumerate}[leftmargin=*]
  \item Compulsory: first time a block is accessed.
  \item Capacity: cache cannot hold all working set blocks.
  \item Conflict: direct-mapped collisions or low associativity.
\end{enumerate}

\textbf{Examples.}
\begin{enumerate}[leftmargin=*]
  \item First reference in a program is compulsory.
  \item Large array that exceeds cache size causes capacity misses.
  \item Two hot arrays mapping to same index cause conflict misses.
\end{enumerate}

\textbf{Common checks and traps.} Conflict misses can be reduced by higher associativity.

\subsection{AMAT and AMSC Formulas (L19, L20)}
\textbf{What it is.} Average Memory Access Time and memory stall cycles quantify cache impact.

\textbf{Step-by-step.}
\begin{enumerate}[leftmargin=*]
  \item AMAT = hit time + (miss rate x miss penalty).
  \item AMSC = (memory requests per instruction) x miss rate x miss penalty.
  \item CPU time = IC x (Ideal CPI + AMSC) x CCT.
\end{enumerate}

\textbf{Examples.}
\begin{enumerate}[leftmargin=*]
  \item Hit 1 cycle, miss 20 cycles, miss rate 5\%: AMAT = 2 cycles.
  \item If 36\% of instructions are loads/stores, D-cache miss adds 1.44 cycles.
  \item CPI goes from 2 to 5.44 in the example from lecture.
\end{enumerate}

\textbf{Common checks and traps.} Use fractions, not percentages, in formulas.

\subsection{Cache Management Policies (L19, L20)}
\textbf{What it is.} Placement, identification, replacement, and write strategy decisions.

\textbf{Step-by-step.}
\begin{enumerate}[leftmargin=*]
  \item Placement: direct-mapped, set-associative, fully associative.
  \item Identification: tag + index + valid bit.
  \item Replacement: LRU, LFU, random.
  \item Writes: write-through or write-back.
\end{enumerate}

\textbf{Examples.}
\begin{enumerate}[leftmargin=*]
  \item Direct-mapped has no replacement choice.
  \item 2-way set-associative chooses one of 2 lines.
  \item Write-back uses a dirty bit per line.
\end{enumerate}

\textbf{Common checks and traps.} Write policies affect miss handling details.

\subsection{Write-Through vs Write-Back and Dirty Bit (L19)}
\textbf{What it is.} Write-through updates memory immediately; write-back delays until eviction.

\textbf{Step-by-step.}
\begin{enumerate}[leftmargin=*]
  \item Write-through: update cache and memory on each store.
  \item Write-back: update cache and set dirty bit.
  \item On eviction, if dirty=1, write line to memory.
\end{enumerate}

\textbf{Examples.}
\begin{enumerate}[leftmargin=*]
  \item Write-through simplifies consistency but increases memory traffic.
  \item Write-back reduces traffic but needs dirty bit.
  \item A write miss in write-back may trigger write of dirty victim.
\end{enumerate}

\textbf{Common checks and traps.} Write-back can cause a write on a read miss.

\subsection{Write-Allocate vs No-Write-Allocate (L19)}
\textbf{What it is.} On a write miss, decide whether to bring the block into cache.

\textbf{Step-by-step.}
\begin{enumerate}[leftmargin=*]
  \item Write-allocate: fetch block, then write into cache.
  \item No-write-allocate: write directly to memory.
  \item Often combined as write-back + write-allocate or write-through + no-write-allocate.
\end{enumerate}

\textbf{Examples.}
\begin{enumerate}[leftmargin=*]
  \item Write-through + no-write-allocate avoids loading on write miss.
  \item Write-back + write-allocate keeps block for future writes.
  \item Write buffer allows CPU to continue while memory writes complete.
\end{enumerate}

\textbf{Common checks and traps.} Use consistent write policy pairs.

\subsection{Set-Associative and Fully Associative Caches (L19, L20)}
\textbf{What it is.} Allow multiple possible cache lines per index to reduce conflicts.

\textbf{Step-by-step.}
\begin{enumerate}[leftmargin=*]
  \item 2-way or n-way set-associative: search all entries in a set.
  \item Fully associative: any block can go anywhere.
  \item More associativity reduces conflict misses but increases hit time.
\end{enumerate}

\textbf{Examples.}
\begin{enumerate}[leftmargin=*]
  \item For 8 lines, direct-mapped = 1-way, 2-way has 4 sets.
  \item Fully associative needs tag compare for all lines.
  \item 2-way set-associative can reduce misses from 5 to 4 in example.
\end{enumerate}

\textbf{Common checks and traps.} Higher associativity needs more comparators.

\begin{figure}[h]
\centering
\fbox{\parbox{0.9\linewidth}{Diagram placeholder: Set-associative cache structure (L19)}}
\end{figure}

\subsection{Replacement Policies (LRU, LFU, Random) (L19)}
\textbf{What it is.} Choose which line to evict on a miss when multiple choices exist.

\textbf{Step-by-step.}
\begin{enumerate}[leftmargin=*]
  \item LRU: evict the least recently used line.
  \item LFU: evict the least frequently used line.
  \item Random: evict a random line, often near-LRU performance.
\end{enumerate}

\textbf{Examples.}
\begin{enumerate}[leftmargin=*]
  \item 2-way cache uses 1 bit to track LRU.
  \item 4-way LRU needs more state and logic.
  \item Random replacement is simplest for high associativity.
\end{enumerate}

\textbf{Common checks and traps.} LRU becomes complex as associativity grows.

\subsection{Cache Performance Example (L19)}
\textbf{What it is.} Compute CPI impact from I-cache and D-cache misses.

\textbf{Step-by-step.}
\begin{enumerate}[leftmargin=*]
  \item Compute I-cache miss cycles: miss rate x penalty.
  \item Compute D-cache miss cycles: load/store fraction x miss rate x penalty.
  \item Add to ideal CPI.
\end{enumerate}

\textbf{Examples.}
\begin{enumerate}[leftmargin=*]
  \item I-miss 2\% with 100-cycle penalty adds 2 cycles.
  \item D-miss 4\% with 36\% loads adds 1.44 cycles.
  \item CPI = 2 + 2 + 1.44 = 5.44.
\end{enumerate}

\textbf{Common checks and traps.} Separate I and D miss contributions.

\subsection{Multiword Blocks, Critical Word First, Early Restart (L20)}
\textbf{What it is.} Larger blocks exploit spatial locality but increase miss penalty; critical word first reduces effective penalty.

\textbf{Step-by-step.}
\begin{enumerate}[leftmargin=*]
  \item On miss, fetch an entire block of multiple words.
  \item Deliver the requested word first (critical word first).
  \item Resume CPU early while the rest of block arrives.
\end{enumerate}

\textbf{Examples.}
\begin{enumerate}[leftmargin=*]
  \item 16-word block reduces compulsory misses for sequential access.
  \item Early restart reduces perceived miss penalty.
  \item Too-large blocks increase conflict and capacity misses.
\end{enumerate}

\textbf{Common checks and traps.} Block size tradeoff depends on workload.

\subsection{Cache Optimization Knobs (L20)}
\textbf{What it is.} Change block size, capacity, associativity, or add levels to improve performance.

\textbf{Step-by-step.}
\begin{enumerate}[leftmargin=*]
  \item Increase block size to reduce compulsory misses.
  \item Increase cache size to reduce capacity misses.
  \item Increase associativity to reduce conflict misses.
  \item Use multi-level caches or software reordering.
\end{enumerate}

\textbf{Examples.}
\begin{enumerate}[leftmargin=*]
  \item Loop interchange improves spatial locality.
  \item L2 cache reduces miss penalty for L1 misses.
  \item Higher associativity increases hit time and power.
\end{enumerate}

\textbf{Common checks and traps.} Improvements can trade one miss type for another.

\subsection{Virtual Memory Basics (L20)}
\textbf{What it is.} MM acts as a cache for disk; programs see a larger virtual address space.

\textbf{Step-by-step.}
\begin{enumerate}[leftmargin=*]
  \item CPU generates logical (virtual) addresses.
  \item OS and hardware translate to physical addresses.
  \item Pages are the VM blocks; page faults occur on misses.
\end{enumerate}

\textbf{Examples.}
\begin{enumerate}[leftmargin=*]
  \item 32-bit virtual address may map to 30-bit physical address.
  \item Each process gets a private virtual address space.
  \item A page fault triggers disk access and OS handler.
\end{enumerate}

\textbf{Common checks and traps.} Page faults are orders of magnitude slower than cache misses.

\subsection{Page Tables and PTE Fields (L20)}
\textbf{What it is.} Page tables map virtual page numbers to physical page numbers.

\textbf{Step-by-step.}
\begin{enumerate}[leftmargin=*]
  \item Split virtual address into page number and offset.
  \item Use page number to index the page table.
  \item PTE stores physical page number and status bits (valid, dirty, reference).
\end{enumerate}

\textbf{Examples.}
\begin{enumerate}[leftmargin=*]
  \item 4 KB pages mean 12-bit offset.
  \item Valid bit 0 indicates page not in memory.
  \item Dirty bit set when a page is modified.
\end{enumerate}

\textbf{Common checks and traps.} Page tables live in physical memory, not in cache by default.

\subsection{Page Faults and Handler Steps (L20)}
\textbf{What it is.} A page fault occurs when the needed page is not in memory.

\textbf{Step-by-step.}
\begin{enumerate}[leftmargin=*]
  \item Trap to OS with faulting address.
  \item Find page on disk; choose victim page.
  \item Write back victim if dirty; load new page.
  \item Update page table and restart instruction.
\end{enumerate}

\textbf{Examples.}
\begin{enumerate}[leftmargin=*]
  \item Page fault costs millions of cycles.
  \item Dirty victim causes extra disk write.
  \item Restart from the faulting instruction after load.
\end{enumerate}

\textbf{Common checks and traps.} Page faults are handled by OS, not hardware alone.

\subsection{Page Replacement and Reference Bits (L20)}
\textbf{What it is.} OS selects a page to evict using algorithms like LRU.

\textbf{Step-by-step.}
\begin{enumerate}[leftmargin=*]
  \item Use reference (use) bit to approximate LRU.
  \item Periodically clear reference bits.
  \item Prefer evicting pages with reference bit 0.
\end{enumerate}

\textbf{Examples.}
\begin{enumerate}[leftmargin=*]
  \item Clock algorithm uses reference bits in a ring.
  \item Dirty page eviction requires write-back.
  \item Clean page eviction avoids disk write.
\end{enumerate}

\textbf{Common checks and traps.} True LRU is expensive to implement in OS.

\subsection{TLB and TLB Miss Handling (L20)}
\textbf{What it is.} TLB caches page table entries to speed address translation.

\textbf{Step-by-step.}
\begin{enumerate}[leftmargin=*]
  \item Check TLB for virtual page number.
  \item On hit, get physical page number and proceed.
  \item On miss, fetch PTE from memory or trigger page fault.
  \item Update TLB and retry.
\end{enumerate}

\textbf{Examples.}
\begin{enumerate}[leftmargin=*]
  \item TLB miss with page in memory: load PTE and retry.
  \item TLB miss with page not in memory: page fault.
  \item Software-managed TLB uses a special exception handler.
\end{enumerate}

\textbf{Common checks and traps.} TLB miss handling can be hardware or software.

\appendix
\section{MIPS Programming Guide (Lecture Notes-Compatible)}

\subsection{Register File and Calling Convention}
\textbf{Key registers.}
\begin{itemize}[leftmargin=*]
  \item $\$zero$ = 0; $\$at$ assembler temp; $\$v0$--$\$v1$ return values.
  \item $\$a0$--$\$a3$ arguments; $\$t0$--$\$t9$ temporaries (caller-saved).
  \item $\$s0$--$\$s7$ saved (callee-saved); $\$sp$ stack pointer; $\$fp$ frame pointer; $\$ra$ return address.
\end{itemize}

\textbf{Example: caller vs callee saved.}
\begin{verbatim}
# caller uses $t0, must save if needed
add $t0, $a0, $a1
sw  $t0, 0($sp)

# callee saves $s0
sw  $s0, 0($sp)
...
lw  $s0, 0($sp)
jr  $ra
\end{verbatim}

\subsection{Data Types, Directives, and Alignment}
\textbf{Common directives.} \texttt{.word}, \texttt{.half}, \texttt{.byte}, \texttt{.asciiz}, \texttt{.space}.

\textbf{Examples.}
\begin{verbatim}
.data
arr:   .word 1, 2, 3, 4
msg:   .asciiz "hello"
buf:   .space 64

.text
la  $t0, arr   # address of array
lw  $t1, 0($t0)
\end{verbatim}

\textbf{Alignment tip.} Word data should be 4-byte aligned for lw/sw.

\subsection{Instruction Formats and Addressing Modes}
\textbf{Formats.} R-type (op, rs, rt, rd, shamt, funct), I-type (op, rs, rt, imm), J-type (op, target).

\textbf{Addressing modes.}
\begin{itemize}[leftmargin=*]
  \item Register: \texttt{add \$t0, \$t1, \$t2}
  \item Immediate: \texttt{addi \$t0, \$t1, 4}
  \item Base+offset: \texttt{lw \$t0, 8(\$sp)}
  \item PC-relative branch: \texttt{beq \$t0, \$t1, label}
\end{itemize}

\subsection{Arithmetic and Logical Operations}
\textbf{Examples.}
\begin{verbatim}
add  $t0, $t1, $t2   # signed add
addu $t0, $t1, $t2   # unsigned add (no overflow trap)
sub  $t0, $t1, $t2
and  $t0, $t1, $t2
or   $t0, $t1, $t2
xor  $t0, $t1, $t2
slt  $t0, $t1, $t2   # $t0=1 if $t1<$t2 (signed)
sltu $t0, $t1, $t2   # unsigned compare
\end{verbatim}

\textbf{Example: absolute value.}
\begin{verbatim}
abs:
  slt  $t0, $a0, $zero
  beq  $t0, $zero, done
  sub  $a0, $zero, $a0
done:
  jr   $ra
\end{verbatim}

\subsection{Shifts and Bit Tricks}
\textbf{Examples.}
\begin{verbatim}
sll $t0, $t1, 2     # multiply by 4
srl $t0, $t1, 1     # logical shift right
sra $t0, $t1, 1     # arithmetic shift right

# extract low 8 bits
andi $t0, $t1, 0xFF
\end{verbatim}

\subsection{Load/Store and Memory Access}
\textbf{Examples.}
\begin{verbatim}
lw  $t0, 0($sp)     # load word
sw  $t0, 4($sp)     # store word
lb  $t1, 0($t0)     # load byte (sign-extend)
sb  $t1, 1($t0)     # store byte
lh  $t2, 0($t0)     # load half
sh  $t2, 2($t0)     # store half
\end{verbatim}

\textbf{Alignment tip.} Use lw/sw on word-aligned addresses only.

\subsection{Branches and Jumps}
\textbf{Examples.}
\begin{verbatim}
beq $t0, $t1, equal
bne $t0, $t1, noteq

loop:
  addi $t0, $t0, -1
  bne  $t0, $zero, loop

j   target
jal func       # jump and link
jr  $ra         # return
\end{verbatim}

\subsection{Functions and Stack Frames}
\textbf{Pattern.}
\begin{verbatim}
func:
  addi $sp, $sp, -16
  sw   $ra, 12($sp)
  sw   $s0, 8($sp)
  move $s0, $a0
  ...
  lw   $s0, 8($sp)
  lw   $ra, 12($sp)
  addi $sp, $sp, 16
  jr   $ra
\end{verbatim}

\textbf{Example: factorial (recursive).}
\begin{verbatim}
fact:
  addi $sp, $sp, -8
  sw   $ra, 4($sp)
  sw   $a0, 0($sp)
  ble  $a0, 1, base
  addi $a0, $a0, -1
  jal  fact
  lw   $a0, 0($sp)
  mul  $v0, $v0, $a0
  j    done
base:
  li   $v0, 1
done:
  lw   $ra, 4($sp)
  addi $sp, $sp, 8
  jr   $ra
\end{verbatim}

\subsection{Syscalls (MARS/QtSPIM)}
\textbf{Common services.}
\begin{itemize}[leftmargin=*]
  \item 1: print int (\texttt{\$a0})
  \item 4: print string (\texttt{\$a0})
  \item 5: read int (returns in \texttt{\$v0})
  \item 10: exit
  \item 2: print float (\texttt{\$f12}), 6: read float (returns in \texttt{\$f0})
\end{itemize}

\textbf{Example: print integer and newline.}
\begin{verbatim}
li   $v0, 1
move $a0, $t0
syscall

li   $v0, 11
li   $a0, 10
syscall
\end{verbatim}

\subsection{Arrays and Loops}
\textbf{Example: sum array of N words.}
\begin{verbatim}
sum:
  move $t0, $a0    # base
  move $t1, $a1    # N
  li   $v0, 0
loop:
  beq  $t1, $zero, done
  lw   $t2, 0($t0)
  add  $v0, $v0, $t2
  addi $t0, $t0, 4
  addi $t1, $t1, -1
  j    loop
done:
  jr   $ra
\end{verbatim}

\textbf{Example: max in array.}
\begin{verbatim}
max:
  lw   $v0, 0($a0)
  addi $a0, $a0, 4
  addi $a1, $a1, -1
loop:
  beq  $a1, $zero, done
  lw   $t0, 0($a0)
  slt  $t1, $v0, $t0
  beq  $t1, $zero, skip
  move $v0, $t0
skip:
  addi $a0, $a0, 4
  addi $a1, $a1, -1
  j    loop
done:
  jr   $ra
\end{verbatim}

\subsection{Strings}
\textbf{Example: strlen.}
\begin{verbatim}
strlen:
  move $t0, $a0
  li   $v0, 0
loop:
  lb   $t1, 0($t0)
  beq  $t1, $zero, done
  addi $v0, $v0, 1
  addi $t0, $t0, 1
  j    loop
done:
  jr   $ra
\end{verbatim}

\subsection{Multiply/Divide Instructions}
\textbf{Examples.}
\begin{verbatim}
mult $t0, $t1
mflo $t2         # low 32 bits of product
mfhi $t3         # high 32 bits of product

div  $t0, $t1
mflo $t2         # quotient
mfhi $t3         # remainder
\end{verbatim}

\subsection{Floating-Point Programming}
\textbf{Example: dot product of 2 floats.}
\begin{verbatim}
dot:
  li   $t0, 0
  mov.s $f0, $f0   # sum in f0
loop:
  beq  $t0, $a2, done
  lwc1 $f2, 0($a0)
  lwc1 $f4, 0($a1)
  mul.s $f6, $f2, $f4
  add.s $f0, $f0, $f6
  addi $a0, $a0, 4
  addi $a1, $a1, 4
  addi $t0, $t0, 1
  j    loop
done:
  jr   $ra
\end{verbatim}

\textbf{Example: int to float conversion.}
\begin{verbatim}
mtc1 $t0, $f0
cvt.s.w $f2, $f0
\end{verbatim}

\subsection{Common Debugging Checklist}
\begin{itemize}[leftmargin=*]
  \item Check register lifetimes: did you overwrite a needed $s$ register?
  \item Check stack alignment and saved $ra$ for nested calls.
  \item Verify load/store offsets and alignment.
  \item Confirm branch logic and delay slots in your environment.
  \item For FP, use correct .s/.d instruction variants.
\end{itemize}

\end{document}
